\scanpage{18}
% L1–4
Caractérisation intrinsèque d'un $n$-cube, déterminé par un ens
\[
\unsure{P|\mathcal P|\Phi}\subset\mathrm{Espaff}(E)
\]
de sous-espaces affines d'un espace affine sur $\unsure{k|h}$ :

% L5–8
a) $P$ \deleted{a un plus petit et un plus grand élément, et}
\inserted{contient $\varnothing$ et $E$ (qui sont donc un plus petit et un plus grand élément), et}
dans $P^{\unsure{**|\times\times}}=P-\{0_P,1_P\}$ ($P$ privé de son plus
petit et plus grand élément) tout élément est
intersection \inserted{(pour la rel. d'ordre)} d'éléments maximaux.

% L9–14
b) Pour tt élément maximal $F$ de $P^{**}$, $\exists!$
él. maximal $F'$ tel que $\mathrm{Inf}_{P^{\deleted{\unsure{**}}}}(F,F')=\varnothing$ [i.e.
$(F,F')$ non minoré dans $P^{**}$ i.e. $F\cap F'$ ne contient
aucun él. de $P$ différent de $\varnothing$] [Ainsi l'ens
$B$ des él maximaux de $P^{**}$ est muni d'une involution $\sigma$ :
$F-F'$ ss pt fixe\unsure{.|..|\ldots{}}]

% L15–21
c) Si à tte $V\in P^{\deleted{\unsure{**}}}$, \inserted{$V\ne\varnothing$}, on associe l'ens \deleted{\unsure{$A$|$B$}} $B_V$
des $b\in B$ qui le majorent, on trouve
une partie \inserted{$A$} de $B$ telle que $b\in A\Longrightarrow\sigma b\notin A$
(i.e. une section partiellement définie de $B$
sur $I=B/\sigma$). Ceci posé, l'application
\[
P\setminus\{\varnothing\}\longrightarrow\Gamma_{\mathrm{part}}(B/I)
\]
ainsi définie, qui est injective par b), est bijective.

% L22–27
\margin{\uline{NB} a) b) c) signifient que l'ens ordonné $P$ est un cube combinatoire}
d) Si $n=\dim\,\mathrm{comb}\,P$, on a \deleted{$E$} $n=\deleted{\unsure{d}}\mathrm{rg}\ E
=\frac12\operatorname{card} B$.

[ a) à d) implique que \inserted{si $V\in P$, $\dim_P V=\dim_E V$, que}
le $\mathrm{Inf}$ \deleted{\unsure{dans $P$|dans $\Phi$}} \inserted{d'une}
partie minorée de $P\setminus\{0\}$ est aussi son
intersection (en tant qu'ens. de parties
de $E$) \inserted{et que les ``hyperplans'' \deleted{\unsure{$E$}} $\in P$ passant par un sommet $s\in P$ sont lin. indép.}. Mais celà n'implique pas que le
% archive sheet number: 17
